\documentclass[11pt]{article}
\usepackage[margin=1.05in]{geometry}
\usepackage[hidelinks]{hyperref}

\title{Stage 3 Proposal: A Constructed Regime-Switching Task\\
\large Vanilla Policy versus Oracle-Signaled Policy on a Task We Control}
\author{Aryan Goyal}
\date{July 22, 2026}

\begin{document}
\maketitle

\noindent
Stage 2 tests the chunk-length switch on natural long-horizon tasks, but it
cannot control how much of each episode is stable transport and how much is
unstable contact, and that mix is exactly what decides how much room the
switch has to win. Stage 3 removes that limitation by constructing the task
ourselves: a repeated transport-and-insert task in robosuite, with $N$ pegs
or sockets spaced far apart on a large table, so that each cycle is a long
free-space transport (open-loop stable, wants a long chunk) followed by a
tight insertion (open-loop unstable, wants per-step feedback), repeated $N$
times. The two quantities the theory cares about become experimental knobs:
the number of regime alternations $N$ (swept over 1, 2, 4) and the ratio of
stable to unstable segment length (set by peg spacing and insertion
tolerance). The prediction is specific: the gap between the oracle-switched
policy and the best fixed chunk length should grow with $N$, because every
added alternation forces any fixed $k$ into one more compromise. The
comparison is vanilla policy at its best fixed $k$ (found by a full sweep,
same protocol as Stage 1a, never a strawman) against the same policy with
oracle switching, and then with predictor switching. Because we define the
task, the regime boundaries are known by construction, which gives three
things the earlier stages had to pay heavily for: a perfect and free oracle
signal (no perturbation labeling), a ground-truth check of the perturbation
labeler (the measured $\lambda$ bands should line up with the designed
segment boundaries), and clean predictor training labels. Demonstrations
come from MimicGen: author the task, record about ten source demonstrations,
and generate one thousand automatically; the estimated cost is one to two
days of task authoring plus a few GPU-days of generation and training per
$N$. The known weakness is that a bespoke task invites the criticism that it
was engineered for the method to win, so Stage 3 is positioned as the
mechanism experiment that explains and quantifies the effect, while Stage 2
remains the naturalistic evidence; the pairing is stronger than either
alone. Stage 3 starts after the Stage 2 results are in, since those results
determine whether it quantifies how the win scales or diagnoses which
regime mix the hypothesis needs.

\end{document}
